\section{Euler's Method}
\label{sec:eulers-method}

\begin{wrapfigure}[]{r}{2in}
  \captionsetup{labelformat=empty}
  \includegraphics*[height=1.8in,width=2in]{../Figures/EulerFig1}
  \label{fig:EulerFig1}
\end{wrapfigure}
Suppose $B$ is a point on the graph of an unknown function $y(x)$ (in
green in the figure at the right). Suppose further that if the line
tangent to the graph at $B$ (in red) crosses the $x$-axis at $(a, 0)$
as shown.
% $x$-axis and the vertical line through $B$ intersect. Find the point %
% $C$ on the $x$-axis where $\overline{CB}$ is tangent to the curve at %
% $B$. The line segment $\overline{CB}$ is the subtangent of the curve
% at  $B$.
% the  slope of the line segment $BC$ is the derivative of $y$ at $B.$
% \centerline{\includegraphics*[height=1in,width=1.5in]{../Figures/subtan2}}
% \begin{wrapfigure}[]{r}{2.3in}
%   \captionsetup{labelformat=empty}
%   \centerline{\includegraphics*[height=1.5in,width=2.3in]{../Figures/subtan1}}
%   \label{fig:subtan1}
% \end{wrapfigure}
% Mathematics in the seventeenth century was still strongly tied to
% Greek geometry. So some of the early, pre-Calculus efforts to find the
% line tangent to a curve at a given point were focused on what was
% called the subtangent problem. 
% That is, 
% 
% is p the subtangent of
% $y$ at $B$ is defined as follows.  Let $A$ be a point where the %
% 
% Note that finding length of $\overline{CA}$ is equivalent to finding
% the derivative of $y$ at $B$ because if $B$ and $C$ are known then
% 
% % \begin{wrapfigure}[]{l}{2in} % % \captionsetup{labelformat=empty}
% %   
% %   \label{fig:subtan2} % % \end{wrapfigure}
% %   \begin{wrapfigure}[]{l}{2in} % % \captionsetup{labelformat=empty}
% %     \includegraphics*[height=1.8in,width=2in]{../Figures/EulerFig1}
% %     \label{fig:EulerFig1} % % \end{wrapfigure} % With the invention of
%     Calculus, the subtangent problem quickly fell % out of favor. However
%     it can be instructive to think about the % construction of tangent
%     lines from this point of view. To be clear, % the subtangent problem
%     is no longer an important problem in % mathematics and we will not be
%     discussing it much further. We only % bring it up because it gives us
%     a nice starting point for the next % topic.  % For the sake of being
%     definite consider the following question:
%     the  length of the subtangent at every point on the
%     curve is equal to $1$, as in the figure at
%     the left?
%     \centerline{\includegraphics*[height=1in,width=2in]{../Figures/subtan2}}

We wish to address the following question\aside{The line segment $AC$
  is called the subtangent of the graph. In the seventeenth mathematics century was still strongly tied to
Greek geometry so some of the early, pre-Calculus efforts to find the
tangent lines were aimed at finding the length of the subtangent of a
curve.
}:
``Is
there a curve which passes through the point $(0, 1)$ with the
property that  the distance from $A$ to $C$ is equal to $1$ regardless
of the location of $B$ on the graph.''

At $B$ form the differential right triangle with the infinitesimal
vertical displacement $\dx{y}$ and the infinitesimal horizontal
displacement $\dx{x}$, as shown. Then the right triangle with sides of
length $y$ and $1$ is proportional to the differential triangle. By
the properties of proportional triangles this curve must satisfy the
differential equation: $  \dfdx{y}{x} = y$.

Since we specified that the curve must pass through the point $(0, 1)$
we see that we need to solve the following IVP.
\begin{equation}
  \label{IVP:ExpDiffeq1}
  \dfdx{y}{x} = y,\  y(0)=1.
\end{equation}
% \begin{wrapfigure}[]{l}{2in}
%   \captionsetup{labelformat=empty}
%   \centerline{\includegraphics*[height=1.5in,width=2in]{../Figures/exp-piecewise-linear}}
%   \label{fig:exp-piecewise-linear}
% \end{wrapfigure}

% Suppose we have a curve with the following property: At every point
% $(x,y)$ we project that tangent line down to the $x$-axis, just as if
% we were implementing Newton's Method. Suppose the tangent line
% intersects the $x$-axis at some point we'll call $\bar{x}$ and that
% the distance between $x$ and $\bar{x}$ is always $1.$ What is the
% shape of such a curve. This situation
% is sketched in the figure to the right.

% As we saw in Section~\ref{sec:witch-agnesi-inverse} knowing the slope
% (derivative) of an unknown function at every point on its graph gives
% us a great deal of information about the unknown function. For example
% we concluded that if $\dfdx{y}{x}=\frac{1}{1+x^2}$ then $y(x)$ is the
% multifunction $\arctan(x)$. By insisting that $y(0)=0$ we were able to
% choose a single branch of the multifunction $\arctan(x)$ which we called $\inverse\tan(x)$.

% The solution of equation~(\ref{eq:ExpDiffeq1}) is also a
% multifunction and  we can  choose a  branch (thus giving us a
% function) by insisting that\aside{You will see why we made this particular
% choice in Section~(\ref{eq:x2diffeq}.)}
% $(0,1)$.  

% Thus we'd like to solve the \term{Initial Value Problem} (IVP)
% \begin{equation}
%   \label{eq:ExpIVP}
%   \dfdx{y}{x} = y,\ y(0)=1.
% \end{equation}
% Notice that the IVP has two parts:
% \begin{enumerate}
% \item  The \term{differential equation}, $
%   \dfdx{y}{x} = y$, and
% \item the \term{initial value}, $y(0)=1$.
% \end{enumerate}


In Chapter~(\ref{cha:last-elem-funct}) we will solve
IVP~(\ref{IVP:ExpDiffeq1}) exactly. For the moment we will be
satisfied if we can find an approximate graph of the solution.  That
will give us a general sense of its shape.


The initial value in IVP~(\ref{IVP:ExpDiffeq1}) shows that  the curve
passes through the point $(0,1)$.  From the differential equation in
IVP~(\ref{IVP:ExpDiffeq1})  we see that
$$
\eval{\dfdx{y}{x}}{x}{0} = y(0) =1.
$$
which means that the curve passes through the point $(0,1)$ with a
slope equal to $1$.  So the equation of the line
tangent to our curve is
$$
y-1=\underbrace{\left[\eval{\dfdx{y}{x}}{x}{0}\right]}_{=y(0)=1}(x-0)
$$
or
$$
y=x+1.
$$
This is not a lot to work with, but let's not set our sights too high.

% \begin{wrapfigure}[]{r}{1.5in}
%   \captionsetup{labelformat=empty}
%   \centerline{\includegraphics*[height=1.5in,width=1.5in]{../Figures/exp-envelope2}}
%   \label{fig:exp-envelope2}
% \end{wrapfigure}
% \begin{wrapfigure}[12]{r}{3in}
%   \vskip-4mm{}
%   \captionsetup{labelformat=empty}
%   \centerline{\includegraphics*[height=1.8in,width=3in]{../Figures/exp-piecewise-linear2}}
%   \label{fig:exp-piecewise-linear2}
% \end{wrapfigure}
By the Principle of Local Linearity, the line tangent to the curve and
the curve itself are going to be nearly indistinguishable near to the
point $(0,1)$. Moreover, we can find any point on the tangent line
since we have its equation.  If we increment the $x$-coordinate just a
little bit while staying on the tangent line then the corresponding
$y-$coordinate on the tangent line will be close to the $y-$coordinate
on the curve. So we can use the $y$-coordinate on the tangent line
(which we know) to approximate the $y$-coordinate on the curve (which
we don't know).

Let's say we increase $x$ by $0.1$ so that $x=0.1$.  The
$y$-coordinate on the line tangent to the curve at $x=0.1$ is $1.1$
(since the equation of our tangent line is $y=x+1$).

We now have two points on our curve: $(0,1)$ and $(0.1, 1.1).$ The
point 
$(0.1, 1.1)$ isn't on the curve, but it's close. Remember  we're
only trying to approximate the curve. Connecting these with a
straight line we have the red segment on the graph below.  So far, so good.
\\
\centerline{\includegraphics*[height=1.8in,width=3in]{../Figures/exp-piecewise-linear2}}

Next we increment $x$ by $0.1$ again, to $x=0.2$.
% It won't
% help to stay on line tangent at $x=0$. That only approximates our curve
% near the point of tangency ($x=0$), and we've moved two steps away from
% it already.
At this point we would really like to have the equation of the line
tangent to the curve at $(0.1, y(0.1))$ but we simply have no way to
obtain it. All we have is the point $(0.1, 1.1)$. But we know that
$y(0.1)\approx1.1$ so the differential equation from
IVP~(\ref{IVP:ExpDiffeq1}) tells us that
$\eval{\dfdx{y}{x}}{x}{0.1}\approx1.1$ also. So we see that the curve will pass
(approximately) through the point $(0.1, 1.1)$ with slope
(approximately) equal to $1.1$. Therefore the equation of the line
tangent to the curve at (approximately) $(0.1, 1.1)$ will be
$y-1.1=1.1(x-0.1)$ and as long as we don't move to far away the
Principle of Local Linearity guarantees that this line and the curve
we seek are close together. So we use the blue line segment from
$(0.1, 1.1)$ to $(0.2, 1.21)$ in our figure above to approximate the
curve on the interval $(0.1, 0.2)$.

If we now simply repeat this process
we compute the points:$(0.3,1.331)$, $(0.4,1.4641)$, $(0.5,1.61051)$, 
and so on.
Plotting  these points and  connecting them with straight line
segments gives us the rest of the sketch above.


% \begin{wrapfigure}[16]{l}{3in}
%   \vskip-4mm{}
%   \captionsetup{labelformat=empty}
%   \centerline{\includegraphics*[height=2.5in,width=3in]{../Figures/exp-piecewise-linear3}}
%   \label{fig:exp-piecewise-linear3}
% \end{wrapfigure}
At each step we use the previous approximation to compute the next, so
it would be miraculous if we actually found exactly the points on the
graph of our curve.  But it should be clear  that the curve we've drawn
will at least resemble the desired curve, as long as we don't
stray too far from our initial point $x=0$.
The sketch below shows our approximation and the actual solution
on the same set of axes\aside{In practice of course, we wouldn't know the
  actual solution -- if we knew that there would be no reason to find an
  approximation -- but in this case we (the authors) do know.}.\\
   \centerline{\includegraphics*[height=2.5in,width=3in]{../Figures/exp-piecewise-linear3}}

% The problem of recreating a curve given its derivative and its
% $y$-coordinate at a single point is called an {\bf Initial Value
% Problem}, or IVP for short.

% An IVP comes in two parts. First there  is the differential equation,
% in our problem:
% $$
% \dfdx{y}{x} =y
% $$
% and second there is the Initial  Condition, in our problem: $y(0)=1.$

The solution of this particular IVP  turns out to be
incredibly useful in mathematics, theoretical physics, engineering,
and science and technology in general.  We will be revisiting it in
the next chapter.

The procedure we have just outlined is known as {\bf Euler's Method};
named for the great eighteenth century mathematician Leonhard
Euler. For now we'll focus on how to use Euler's Method to approximate the
solution of an arbitrary IVP. We formally define and Initial Value
Problem as follows.

\begin{mydefinition}[Initial Value Problem]
  \label{def:IVP}
  An \term{Initial Value Problem} is a differential equation of the
  form $$\dfdx{y}{x}=f(x,y),$$ along with an \term{Initial Value} $$y(x_0)=y_0.$$
\end{mydefinition}

Don't let the formalism of this definition scare you. All it says is
that at every point $(x,y)$ the slope of $y(x)$,
$\left(\dfdx{y}{x}\right)$, is given by some formula, $f(x,y)$,
which may involve both $x$ and $y$\aside{In our problem above we had
  $\dfdx{y}{x}=f(x,y)=y$.}, and that we know the value of $y(x)$ for a 
single value of $x$. Specifically, at  $x=x_0$.

To write down Euler's Method clearly we will need some
notation.  % Assume that $y$ is the unknown function of $x$ that we
% are trying to find and  denote it in the usual way: $y=y(x).$
% We let $f(x,y)$ denote some particular (but
% unspecified) formula\aside{In our example $f(x,y)$ was
% particularly simple: $f(x,y)=y.$ } involving one or both of $x$ and $y.$ Since
% $\dfdx{y}{x}=f(x,y)$ we see that the slope of the graph of $y$ can
% by found if we know both $x$ and $y$ by simply evaluating $f(x,y).$

% So  if we know that the graph of $y(x)$ passes through some
% specific point, say $(x_0,y_0)$ then it will pass through with
% slope: $\left.\dfdx{y}{x}\right|_{(x_0,y_0)}=f(x_0,y_0).$
% The point $(x_0,y_0)$  the  \term{Initial Value} of
% $y(x)$: $y(x_0)=y_0.$


% Don't let all of this notation throw you off. We just need a way to
% talk about a generic, rather than some specific, IVP.

% In general an IVP takes the following form:
% $$
% \dfdx{y}{x} = f(x,y),\ \ \ y(x_0)=y_0.
% $$
% This says that we can find the slope of $y(x)$ at any point $(x,y)$ by
% computing $f(x,y),$ and that the point $(x_0,y_0)$ is on the graph of
% $y(x).$ In fact, $(x_0, y_0)$ is the only point we know for certain is on the graph of $y(x)$. That's why it's call the Initial Value.

We know that the equation of the tangent line to the curve $y=y(x)$ at
the point $(x_0,y_0)$ is given by 
\begin{align*}
  % y=y_0+\underbrace{\left[\left.\dfdx{y}{x}\right|_{(x_0,y_0)}\right]}_{=f(x_0,y_0)}(x-x_0)=y_0+f(x_0,y_0)(x-x_0
  % ).
  y\amp{}=y_0+\left(\eval{\dfdx{y}{x}}{(x,y)}{(x_0,y_0)}\right)(x-x_0)\\
  \amp{}=y_0+f(x_0,y_0)(x-x_0 ).
\end{align*}
So if we choose $x_1$ very close to $x_0$ and compute
$
y_1=y_0+f(x_0,y_0 )(x_1-x_0 )
$
then $(x_1,y_1)$ would be approximately on the curve. In this way we can generate
a sequence of points
$(x_0,y_0 ), (x_1,y_1 ), (x_2,y_2 ), (x_3,y_3 ),\ldots$ which are
approximately on the curve.  Connecting them with straight line
segments should provide an
approximate graph of the curve. 

\noindent{}Formalizing all of this we have:\\
\vskip2mm{}
\hrule{}
\vskip2mm{}
\centerline{\bf \Large \underline{Euler's Method}}
\vskip3mm{}
\noindent{}Given an IVP:
$$
\dfdx{y}{x} = f(x,y),\ \ \ y(x_0)=y_0.
$$
approximate points $(x_1,y_1 ), (x_2,y_2 ), (x_3,y_3 ),\ldots$   on
the graph  of $y(x)$ by computing:
$$
y_{n+1}=f(x_n,y_n )(x_{n+1}-x_n)+y_n
\text{ for } n=1, 2, 3, \ldots.$$
% \vskip2mm
\hrule{}
\vskip6mm

At each iteration the value of $y$ has been approximated, so the next
approximation is probably not as good. Thus, as we move further from
our initial value $(x_0, y_0)$ our approximation probably deviates
further away from the actual curve.  However,  near to the
initial value, $(x_0,y_0)$ we should have reasonable approximation
to the curve $y=y(x)$. The next two
problems demonstrate this. 

\begin{embeddedproblem}{}
  \label{embed:EulerSine}
  Observe that $y=\sin(x)$ satisfies the IVP
  $$
  \dfdx{y}{x}=\cos(x),\ \ \  y(0)=0.
  $$

  Use Euler's Method on this IVP to complete the table. Then plot the
  points you generated and the graph of $y=\sin(x)$ on the same
  set of axes so you can compare them.\\\\
  \centerline{\includegraphics*[height=1.4in,width=5.1in]{../Figures/EulerTable1}}
  % \vskip-3mm{}
\end{embeddedproblem}
%%%% Warning!!! Don't delete the commented tables below. They provide
%%%% the images that we're using for these tables. -- ecb5  -- Thu Feb 28 18:20:04 2019 -- 
% $$
% \begin{array}{|c|c|c||c|c|c|}
%   \hline{}
%   \bf n\amp{}\bf x_n\amp{}\bf y_n=y_{n-1}+\cos(x_{n-1})\cdot(x_n-x_{n-1})\amp{}\bf n\amp{}\bf
%   x_n\amp{}\bf
%   y_n=y_{n-1}+\cos(x_{n-1})\cdot(x_n-x_{n-1})\\\hline
    %     \bf  0\amp{}\bf 0  \amp{}\amp{}\bf    8\amp{}\bf 0.8\amp{}\\\hline
%   \bf  1\amp{}\bf 0.1\amp{}\amp{}\bf    9\amp{}\bf 0.9\amp{}\\\hline
%   \bf  2\amp{}\bf 0.2\amp{}\amp{}\bf  10\amp{}\bf 1.0\amp{}\\\hline
%   \bf  3\amp{}\bf 0.3\amp{}\amp{}\bf 11\amp{}\bf 1.1 \amp{}\\\hline
%   \bf  4\amp{}\bf 0.4\amp{}\amp{}\bf 12\amp{}\bf 1.2 \amp{}\\\hline
%   \bf  5\amp{}\bf 0.5\amp{}\amp{}\bf 13\amp{}\bf 1.3 \amp{}\\\hline
%   \bf  6\amp{}\bf 0.6\amp{}\amp{}\bf 14\amp{}\bf 1.4 \amp{}\\\hline
%   \bf  7\amp{}\bf 0.7\amp{}\amp{}\bf 15\amp{}\bf 1.5 \amp{}\\\hline
% \end{array}
% $$
% \vskip-6mm{}


\begin{embeddedproblem}{}
  \label{embed:EulerCosine}
  Observe that $y=\cos(x)$ satisfies the IVP
  $$
  \dfdx{y}{x}=-\sin(x),\ \ \  y(0)=1.
  $$

  Use Euler's Method on this IVP to complete the following table.
  Then plot the points you generated and the graph of $y=\cos(x)$ on
  the same
  set of axes so you can compare them.\\\\
  \centerline{\includegraphics*[height=1.5in,width=5.1in]{../Figures/EulerTable2}}
%  \vskip-3mm{}
  %%%% Warning!!! Don't delete the commented tables below. They provide
  %%%% the images that we're using for these tables. -- ecb5  -- Thu Feb 28 18:20:04 2019 -- 
  % $$
  % \begin{array}{|c|c|c||c|c|c|}
  %   \hline{}
  %   n\amp{}\bf x_n\amp{}\bf y_n=y_{n-1}-\sin(x_{n-1})\cdot(x_n-x_{n-1})\amp{}\bf   n\amp{}\bf x_n\amp{}\bf y_n=y_{n-1}-\sin(x_{n-1})\cdot(x_n-x_{n-1})\\\hline
  %   \bf 0\amp{}\bf 0  \amp{}\bf \amp{}\bf     8\amp{}\bf 0.8\amp{}\bf \\\hline
  %   \bf 1\amp{}\bf 0.1\amp{}\bf \amp{}\bf     9\amp{}\bf 0.9\amp{}\bf \\\hline
  %   \bf 2\amp{}\bf 0.2\amp{}\bf \amp{}\bf   10\amp{}\bf 1.0\amp{}\bf \\\hline
  %   \bf 3\amp{}\bf 0.3\amp{}\bf \amp{}\bf 11\amp{}\bf 1.1\amp{}\bf \\\hline
  %   \bf 4\amp{}\bf 0.4\amp{}\bf \amp{}\bf 12\amp{}\bf 1.2\amp{}\bf \\\hline
  %   \bf 5\amp{}\bf 0.5\amp{}\bf \amp{}\bf 13\amp{}\bf 1.3\amp{}\bf \\\hline
  %   \bf 6\amp{}\bf 0.6\amp{}\bf \amp{}\bf 14\amp{}\bf 1.4\amp{}\bf \\\hline
  %   \bf 7\amp{}\bf 0.7\amp{}\bf \amp{}\bf 15\amp{}\bf 1.5\amp{}\bf \\\hline
  % \end{array}
  % $$
\end{embeddedproblem}

%\pagebreak[4]
  \begin{embeddedproblem}
    Use Euler's Method to approximate the solutions of the given IVPs
    by constructing a table like the ones in
    Problems~(\ref{embed:EulerSine})
    and~(\ref{embed:EulerCosine}). Then plot the points you found,
    connecting them with straight line segments.
    \begin{enumerate}[label={\bf{}(\alph*)}]
      \begin{multicols}{2}
      \item $\dfdx{y}{x}=y^2,\ \ \ y(0)=1$
      \item $\dfdx{y}{x}=y^3,\ \ \ y(0)=1$
      \item $\dfdx{y}{x}=y^4,\ \ \ y(0)=1$
      \item $\dfdx{y}{x}=y^5,\ \ \ y(0)=1$
     \end{multicols}
    \end{enumerate}
  \end{embeddedproblem}

%  \vskip-5mm{}
  \begin{embeddedproblem}
    Use Euler's Method to approximate the solutions of the given IVPs
    by constructing a table like the ones in
    Problems~(\ref{embed:EulerSine})
    and~(\ref{embed:EulerCosine}). Then plot the points you found,
    connecting them with straight line segments.
    \begin{enumerate}[label={\bf{}(\alph*)}]
      \begin{multicols}{2}
      \item $\dfdx{y}{x}=\frac{1}{y},\ \ \ y(1)=1$
      \item $\dfdx{y}{x}=\frac{1}{y^2},\ \ \ y(1)=1$
      \item $\dfdx{y}{x}=\frac{1}{y^3},\ \ \ y(1)=1$
      \item $\dfdx{y}{x}=\frac{1}{y^4},\ \ \ y(1)=1$ 
      \end{multicols}
    \end{enumerate}
  \end{embeddedproblem}

%\pagebreak[4]
  \begin{embeddedproblem}{}
    Use Euler's Method to approximate the solutions of the given IVPs
    by constructing a table like the ones in
    Problems~(\ref{embed:EulerSine})
    and~(\ref{embed:EulerCosine}). Then plot the points you found,
    connecting them with straight line segments.
    \begin{enumerate}[label={\bf{}(\alph*)}]
      \begin{multicols}{3}
      \item $\dfdx{y}{x}=xy,\ \ \ y(0)=1$
      \item $\dfdx{y}{x}=x^2y,\ \ \ y(0)=1$
      \item $\dfdx{y}{x}=xy^2,\ \ \ y(0)=1$
      \item $\dfdx{y}{x}=x^2y^2,\ \ \ y(0)=1$ 
      \item $\dfdx{y}{x}=\frac{x}{y},\ \ \ y(1)=1$
      \item $\dfdx{y}{x}=\frac{x^2}{y},\ \ \ y(1)=1$
      \item $\dfdx{y}{x}=\frac{x}{y^2},\ \ \ y(1)=1$
      \item $\dfdx{y}{x}=\frac{x^2}{y^2},\ \ \ y(1)=1$
      \item $\dfdx{y}{x}=\frac{x^2}{y^2},\ \ \ y(1)=2$
      \end{multicols}
    \end{enumerate}
  \end{embeddedproblem}    

 %  For now, we will be content with approximations.   Later, we will
 %  learn techniques that will allow us to solve some IVPs exactly.
 % Let's apply
 %  Euler's Method to a couple of physical problems.
  
%\pagebreak[4]
\begin{embeddedproblem}{}
  \label{prob:pursuit}
  Suppose a rocket $P$ starts at the point $(1,0)$ and travels up
  the line $x=1$ at a constant speed $v.$ When it is at $(1,0),$ a
  homing missile $M$ is fired from the origin directly at the
  rocket. We want to find the path the missile will follow, assuming
  that it travels at a speed which is $k$ times the speed of the rocket
  $(k\gt{}1)$ and is always aimed directly at the rocket.  The diagram
  below shows the situation at time $t.$\\
  \centerline{\includegraphics*[height=2in,width=3in]{../Figures/PursuitCurve}}
  \begin{enumerate}[label={\bf{}(\alph*)}]
  \item If we let $s$ denote the distance the missile has traveled at
    time $t,$ show that the missile's path must satisfy the IVP
    $$
    \dfdx{y}{x}=\frac{\frac{s}{k}-y}{1-x},\ \       y(0)=0.
    $$
    \hint{Find a triangle which is proportional to the differential
      triangle shown.}
  \item Setting $k=1.5,$ apply Euler's Method to fill in the following
    table.\\
    \comment{There is a little wrinkle here. You'll need to keep track
    of  the length of $s$ which we have no way to compute. However, you can
    approximate the length of $s$ with the length of the polygonal
    path you've created up to that point.  This is why the extra column for values of
    $s$ is included.}\\
    \leftline{\includegraphics*[height=1.25in,width=5in]{../Figures/EulerTable4}}
    
    % Warning!!! Don't delete the commented tables below. They provide
    % the images that we're using for these tables. -- ecb5  -- Thu Feb 28 18:20:04 2019 -- 
    % $$
    % \begin{array}{|c|c|c|c|c||c|c|c|}
    %   \hline{}
    %   \bf{}  n\amp{}\bf x_n\amp{}\bf y_n=y_{n-1}+\left(\left.\dfdx{y}{x}\right|_{x_{n-1}}\right)\cdot(x_n-x_{n-1})\amp{}\bf s\amp{}\bf   n\amp{}\bf x_n\amp{}\bf y_n=y_{n-1}+\left(\left.\dfdx{y}{x}\right|_{x_{n-1}}\right)\cdot(x_n-x_{n-1})\amp{}\bf s\\\hline
    %   \bf{} 0\amp{}\bf 0  \amp{}\bf 0\amp{}\bf 0\amp{}\bf     5\amp{}\bf 0.5\amp{}\bf \amp{}\bf \\\hline
    %   \bf{}  1\amp{}\bf 0.1\amp{}\bf  \amp{}\bf \amp{}\bf     6\amp{}\bf 0.6\amp{}\bf \amp{}\bf \\\hline
    %   \bf{}  2\amp{}\bf 0.2\amp{}\bf \amp{}\bf \amp{}\bf     7\amp{}\bf 0.7\amp{}\bf \amp{}\bf \\\hline
    %   \bf{}  3\amp{}\bf 0.3\amp{}\bf \amp{}\bf \amp{}\bf     8\amp{}\bf 0.8\amp{}\bf \amp{}\bf \\\hline
    %   \bf{}  4\amp{}\bf 0.4\amp{}\bf \amp{}\bf \amp{}\bf     9\amp{}\bf 0.9\amp{}\bf \amp{}\bf \\\hline
    % \end{array}
    % $$
  \item Plot the points obtained in part b to see if this resembles the
    path you think the missile would follow.
  \end{enumerate}
\end{embeddedproblem}

\begin{embeddedproblem}{}
  \label{prob:Tractrix}
  Consider the top view of a tractor-trailer as it turns, as shown
  below.  Initially, the center of the rear axle of the tractor is at
  the origin and the center of
  the rear axle of the trailer is at the point $(1,0).$\\
  \centerline{
%    \includegraphics*[height=1.5in,width=2.5in]{../Figures/Tractrix}}
    \includegraphics*[height=2.3in,width=3.5in]{../Figures/Tractrix}}
  \begin{enumerate}[label={\bf{}(\alph*)}]
  \item Suppose the tractor pulls the front wheels up the $y$-axis and
    that the rear wheels don't slip.  Show that the path of 
    the center of the rear axle of the trailer follows, $y=y(x)$  must satisfy the
    Initial Value Problem
    \begin{equation}
      \label{eq:TractrixIVP}
      \dfdx{y}{x}=-\frac{\sqrt{1-x^2 }}{x}\ \ y(1)=0.
  \end{equation}
    \hint{Find a differential triangle and an ordinary triangle which
      are proportional like in Problem~\#\ref{prob:pursuit}.}
  \item This path you found in part (b) is called a tractrix from the
    Latin verb \foreign{trahere}, meaning ``to drag
    or pull.'' 
    Apply Euler's Method to complete the following table to
    approximate the tractrix.\\
    \leftline{\includegraphics*[height=1.25in,width=5in]{../Figures/EulerTable3}}
    % Warning!!! Don't delete the commented tables below. They provide
    % the images that we're using for these tables. -- ecb5  -- Thu Feb 28 18:20:04 2019 -- 
    % $$
    % \begin{array}{|c|c|c||c|c|c|}
    %   \hline{}
    %   \bf  n\amp{}\bf x_n\amp{}\bf y_n=y_{n-1}+\left(\left.\dfdx{y}{x}\right|_{x_{n-1}}\right)\cdot(x_n-x_{n-1})\amp{}\bf   n\amp{}\bf x_n\amp{}\bf y_n=y_{n-1}+\left(\left.\dfdx{y}{x}\right|_{x_{n-1}}\right)\cdot(x_n-x_{n-1})\\\hline
    %   \bf 0\amp{}\bf 1  \amp{}\bf \amp{}\bf     5\amp{}\bf 0.5\amp{}\bf \\\hline
    %   \bf 1\amp{}\bf 0.9\amp{}\bf \amp{}\bf     6\amp{}\bf 0.4\amp{}\bf \\\hline
    %   \bf 2\amp{}\bf 0.8\amp{}\bf \amp{}\bf     7\amp{}\bf 0.3\amp{}\bf \\\hline
    %   \bf 3\amp{}\bf 0.7\amp{}\bf \amp{}\bf     8\amp{}\bf 0.2\amp{}\bf \\\hline
    %   \bf 4\amp{}\bf 0.6\amp{}\bf \amp{}\bf     9\amp{}\bf 0.1\amp{}\bf \\\hline
    % \end{array}
    % $$
  \item  Plot the points obtained in part b to see if this looks like
    the path the rear axle of the trailer would take.
  \end{enumerate}

  We will re-visit this problem in Section~\ref{subsec:tractrix}.
\end{embeddedproblem}





%%% Local Variables: 
%%% outline-minor-mode: t
%%% eval:(flyspell-mode)
%%% TeX-master: "DiffCalc"
%%% End: 
